% EMNLP 2026 System Demonstrations submission for Juddges App.
%
% Single-blind: author names ARE included at submission time.
% Page limit: 6 pages of content + unlimited references/ethics/appendix.
% Accepted papers receive 1 bonus page.
%
% This skeleton uses the official ACL style. See README.md for how to
% fetch acl.sty and acl_natbib.bst from the acl-org repository.

\documentclass[11pt]{article}

\usepackage[]{acl}

\usepackage{times}
\usepackage{latexsym}
\usepackage[T1]{fontenc}
\usepackage[utf8]{inputenc}
\usepackage{microtype}
\usepackage{graphicx}
\usepackage{booktabs}
\usepackage{multirow}
\usepackage{url}
\usepackage{hyperref}
\usepackage{xcolor}

\newcommand{\todo}[1]{\textcolor{red}{\textbf{[TODO: #1]}}}

\title{Juddges App: An Open-Source Web Platform for Hybrid Search,
Retrieval-Augmented Chat, and Structured Extraction over
Polish and England \& Wales Court Judgments}

\author{
  \L{}ukasz Augustyniak\textsuperscript{1} \quad
  Jakub Binkowski\textsuperscript{1} \quad
  Albert Sawczyn\textsuperscript{1} \\
  \textbf{Kamil Tagowski\textsuperscript{1}} \quad
  \textbf{Micha\l{} Bernaczyk\textsuperscript{2}} \quad
  \textbf{Krzysztof Kami\'{n}ski\textsuperscript{3}} \quad
  \textbf{Tomasz Kajdanowicz\textsuperscript{1}} \\
  \textsuperscript{1}Wroc\l{}aw University of Science and Technology, Poland \\
  \textsuperscript{2}University of Wroc\l{}aw, Poland \\
  \textsuperscript{3}Court of Appeal, Wroc\l{}aw, Poland \\
  \texttt{lukasz.augustyniak@pwr.edu.pl}
}

\begin{document}
\maketitle

\begin{abstract}
We present \textbf{Juddges App}, an open-source, self-hostable web
platform for searching, analysing, and extracting structured information
from judicial decisions of Polish appellate courts and the Court of
Appeal of England \& Wales. The system combines hybrid retrieval (dense
embeddings in \texttt{pgvector} with full-text indexing in Meilisearch),
retrieval-augmented chat grounded in matched cases, schema-driven
extraction via a LangGraph agent, an inline annotation workspace, and
a corpus-level analytics dashboard. Sample corpora are loaded on demand
from two public Hugging Face datasets so that any third party can
reproduce the stack against the same data used by the maintainers, or
substitute their own legal corpus. We report retrieval and extraction
results on \todo{N} held-out queries and document an end-to-end demo
walkthrough. The platform, datasets, and benchmarks are released under
Apache~2.0 at \url{https://github.com/pwr-ai/juddges-app}.
\end{abstract}

\section{Introduction}
\label{sec:intro}

Empirical legal research and the day-to-day work of legal practitioners
both depend on the ability to locate, compare, and synthesise relevant
court decisions across large case-law corpora
\citep{katz2023legaltech,chalkidis2020legalbert}. Existing commercial
legal-search portals are closed, jurisdiction-bound, often paywalled,
and rarely expose retrieval signals or underlying models, which hinders
reproducible research and prevents fine-grained extensions such as
custom extraction schemas or jurisdiction-specific embeddings. Open data
initiatives have begun to address the corpus-side gap
\citep{chalkidis2022lexglue,niklaus2023multilegalpile,guha2023legalbench},
but a research-grade, deployable application that ties together hybrid
retrieval, retrieval-augmented question answering, and structured
extraction over multilingual case law has been missing.

\textbf{Juddges App} is designed to close that gap. It is the deployable
companion to the JuDDGES research initiative \citep{juddges-project},
which curates the underlying corpora and develops the legal-NLP
methodology. Both projects are released independently under Apache~2.0.
This demonstration paper makes three contributions:

\begin{itemize}
  \item A transparent, end-to-end open-source reference stack covering
    hybrid retrieval, RAG-grounded chat, and schema-driven extraction
    for multilingual case law.
  \item A reproducible evaluation protocol with released queries and
    held-out gold annotations for retrieval and extraction.
  \item A self-hostable platform that has been used to investigate
    cross-jurisdictional case comparison between Polish and
    England~\&~Wales appellate judgments.
\end{itemize}

\section{Related Work}
\label{sec:related}

\paragraph{Legal IR and case-law search.}
\todo{2--3 sentences contrasting Juddges App with prior legal-IR
demos and platforms (e.g., CaseLaw access projects, Lexlytica,
open-data portals). Cite recent ACL/EMNLP/NAACL demo papers on legal
search.}

\paragraph{Retrieval-augmented LLMs for law.}
\todo{2--3 sentences on RAG-over-case-law work and how the chat chain
here differs (grounded citations, hybrid retriever, Polish + English).}

\paragraph{Information extraction from judgments.}
\todo{2--3 sentences on schema-driven extraction from legal text and
how the LangGraph agent here generalises beyond fixed templates.}

\section{System Description}
\label{sec:system}

\subsection{Architecture}
\label{sec:architecture}

\begin{figure*}[t]
  \centering
  \includegraphics[width=0.95\textwidth]{figures/architecture.pdf}
  \caption{\textbf{Juddges App architecture.} A Next.js frontend talks
  to a FastAPI backend that exposes LangChain chains via LangServe.
  Judgments are indexed both in a \texttt{pgvector} HNSW table
  (semantic) and in Meilisearch (lexical); the hybrid retriever
  combines both signals. Celery handles asynchronous extraction jobs.
  \todo{Generate architecture diagram.}}
  \label{fig:architecture}
\end{figure*}

The system is implemented as a TypeScript/Python monorepo: a Next.js~15
\citep{nextjs} frontend, a FastAPI \citep{fastapi} backend that exposes
LangChain \citep{langchain} chains via LangServe, a Supabase-managed
PostgreSQL database with \texttt{pgvector} \citep{pgvector} for semantic
indexes, and Meilisearch \citep{meilisearch} for full-text search.
Figure~\ref{fig:architecture} summarises the data flow. The backend
ships as three reusable Poetry packages---\texttt{juddges\_search},
\texttt{schema\_generator\_agent}, and \texttt{research\_agent}---each
documented so that retrieval, extraction, and agent components can be
embedded into other legal-NLP applications without the web frontend.

\subsection{Hybrid Retrieval}
\label{sec:retrieval}

Each judgment is indexed both as a \texttt{vector(768)} embedding in a
\texttt{pgvector} HNSW index and as a full-text record in Meilisearch.
A configurable reranker combines the two signals at query time so that
lexical-precise queries (citations, statute references) and semantic
queries (factual analogies, cross-lingual paraphrases) share a single
search surface. \todo{Describe the rerank formula or the score
combination strategy actually used.}

\subsection{Retrieval-Augmented Chat}
\label{sec:chat}

A LangChain chat chain retrieves top-$k$ judgments and grounds the
response in matched passages, returning explicit citations back to the
original cases. \todo{Specify the prompt template, citation extraction
method, and refusal behaviour for low-confidence answers.}

\subsection{Schema-Driven Extraction}
\label{sec:extraction}

A LangGraph agent generates a Pydantic schema from a natural-language
description of the target fields, runs LLM extraction against a chosen
corpus slice, and persists results as a versioned batch for downstream
analysis. \todo{Describe the schema-generation prompt, the validation
loop, and how repair is triggered on malformed outputs.}

\subsection{Annotation Workspace and Analytics}
\label{sec:annotation}

A TipTap-based rich-text editor lets users annotate judgments inline
and persists annotations server-side, so the same corpus can serve
research analyses and human-in-the-loop training-data construction.
The analytics dashboard renders pre-computed aggregations over
jurisdiction, court, and decision-trend slices.

\section{Evaluation}
\label{sec:eval}

We evaluate the three components that are most consequential for
research use: retrieval, extraction, and end-to-end RAG groundedness.

\subsection{Retrieval}
\label{sec:eval-retrieval}

\paragraph{Setup.} We evaluate on \todo{N} queries split across
Polish (\texttt{sample\_queries.json}) and cross-jurisdictional
(\texttt{cross\_jurisdictional\_queries.json}) sets. Relevance judgments
are \todo{describe annotation protocol and inter-annotator agreement
$\kappa$}.

\paragraph{Metrics.} Recall@$k$ for $k \in \{1, 5, 10\}$, MRR, and
nDCG@10.

\paragraph{Results.}
Table~\ref{tab:retrieval} shows that hybrid retrieval outperforms
either signal alone.

\begin{table}[t]
  \centering
  \small
  \begin{tabular}{lcccc}
    \toprule
    Retriever & R@1 & R@5 & R@10 & MRR \\
    \midrule
    BM25 (Meilisearch) & \todo{} & \todo{} & \todo{} & \todo{} \\
    Dense (pgvector)   & \todo{} & \todo{} & \todo{} & \todo{} \\
    Hybrid             & \todo{} & \todo{} & \todo{} & \todo{} \\
    \bottomrule
  \end{tabular}
  \caption{Retrieval performance on \todo{N} held-out queries.}
  \label{tab:retrieval}
\end{table}

\subsection{Information Extraction}
\label{sec:eval-extraction}

\paragraph{Setup.} A gold set of \todo{N} judgments annotated with
\todo{describe schema fields, e.g., court, judges, charges, verdict,
sentence}.

\paragraph{Metrics.} Exact-match accuracy and span F1 per field;
macro-averaged across fields. \todo{Report per-field numbers in
Table~\ref{tab:extraction} (appendix).}

\subsection{End-to-End RAG}
\label{sec:eval-rag}

We report groundedness rate (fraction of answers fully supported by
retrieved passages) and citation accuracy on a \todo{N}-question
evaluation set sampled from real user queries.

\subsection{Latency}
\label{sec:eval-latency}

\paragraph{Setup.} Single-tenant production deployment,
\todo{describe hardware}.

\paragraph{Results.} Median end-to-end latency: \todo{}~ms search,
\todo{}~ms first chat token, \todo{}~ms per-judgment extraction.

\section{Demo Walkthrough}
\label{sec:demo}

\begin{figure}[t]
  \centering
  \includegraphics[width=0.95\columnwidth]{figures/search-ui.pdf}
  \caption{Hybrid search results with side-by-side excerpt previews
  and citation-ready exports. \todo{Add annotated screenshot.}}
  \label{fig:search-ui}
\end{figure}

The accompanying 2.5-minute screencast (linked in the submission)
demonstrates four flows: (1) hybrid search across the Polish and
England \& Wales corpora, including cross-lingual queries
(Section~\ref{sec:retrieval}); (2) grounded chat answering a legal
question with inline citations (Section~\ref{sec:chat}); (3) building
a Pydantic schema from a natural-language description and running it
over a corpus slice (Section~\ref{sec:extraction}); and (4) inline
annotation that is persisted server-side and surfaced in analytics
(Section~\ref{sec:annotation}). Figure~\ref{fig:search-ui} shows the
search interface; additional screenshots are in
Appendix~\ref{app:screenshots}.

\section{Conclusion}
\label{sec:conclusion}

We presented Juddges App, an Apache-2.0 web platform that combines
hybrid retrieval, retrieval-augmented chat, schema-driven extraction,
and annotation over Polish and England~\&~Wales court judgments. The
system, datasets, and benchmarks are released to support reproducible
legal-NLP research and self-hosted deployments by legal practitioners.

\section*{Ethical Considerations}
\label{sec:ethics}

Court judgments may contain personal data of parties, witnesses, and
victims. We rely on the publicly released JuDDGES datasets, which
follow the source courts' publication policies. Polish appellate
judgments are published by the courts after redaction of sensitive
personal data; we do not perform additional redaction and we surface
the original publisher and case identifier with every result so users
can verify provenance. We do not store identifying information about
end users beyond what Supabase Auth requires for authentication.
Retrieval and chat outputs reflect the biases of the underlying
embeddings and language models; the application is therefore positioned
as a research and triage tool, \emph{not} a substitute for legal
counsel or judicial decision-making. \todo{Expand if Industry-track
reviewers flag this; for Demo it is sufficient.}

\section*{Limitations}
\label{sec:limitations}

\todo{Optional for Demo track, but recommended. Cover: (i) jurisdiction
coverage limited to PL appellate criminal and EN\&W appellate; (ii)
embedding model and chat backbone choices may not generalise to other
languages; (iii) extraction quality depends on schema specification
and underlying LLM; (iv) no formal user study yet.}

\section*{Acknowledgements}

We thank the JuDDGES research project at Wroc\l{}aw University of
Science and Technology for the curated Polish and England~\&~Wales
judgment corpora and the legal-NLP methodology this application builds
on, and the Court of Appeal in Wroc\l{}aw for domain expertise on
Polish appellate procedure. \todo{Add explicit grant numbers if
applicable.}

\bibliography{paper}
\bibliographystyle{acl_natbib}

\appendix

\section{Additional Screenshots}
\label{app:screenshots}

\todo{Add full-size screenshots of chat, extraction-batch dashboard,
and analytics views.}

\section{Reproducibility Details}
\label{app:repro}

Code: \url{https://github.com/pwr-ai/juddges-app}. Persistent DOI:
\url{https://doi.org/10.5281/zenodo.19911856}. Datasets: PL
(\url{https://huggingface.co/datasets/JuDDGES/pl-appealcourt-criminal},
DOI 10.57967/hf/8772) and EN
(\url{https://huggingface.co/datasets/JuDDGES/en-appealcourt}, DOI
10.57967/hf/8773). Models used in reported numbers: \todo{embedding
model name and version, chat backbone name and version, extraction
backbone name and version}. Hardware: \todo{}.

\end{document}
